\section{2025-09-18 - Beugung}

\startitemize[packed]
\item
  Interferenz ist auch an einem \quotation{einzel Spalt} möglich
\stopitemize

\startbuffer[fig:einzelspalt]
\startMPpage

u := 1cm;

draw (0, 0.5 * u)--(1.5 * u, 0.5 * u);
draw (0, -0.5 * u)--(1.5 * u, -0.5 * u);

draw (1.5 * u, 0.5 * u)--(1.5 * u, -0.5 * u);
draw (2.5 * u, 0.5 * u)--(2.5 * u, -0.5 * u);

draw (4 * u, 0.5 * u)--(2.5 * u, 0.5 * u);
draw (4 * u, -0.5 * u)--(2.5 * u, -0.5 * u);

drawarrow (2 * u, -1.5 * u)--(1.5 * u, -0.6 * u);

drawarrow (2 * u, -1.5 * u)--(2.5 * u, -0.6 * u);

% label.bot("Kante gilt als 'Spalt'", (0.5 * u, -1.5 * u));
label.bot("Kanten gelten als 'Spalt'", (2 * u, -1.5 * u));

\stopMPpage

\stopbuffer
\placefigure[here][fig:einzelspalt]{Beugung am Einzelspalt}{
\typesetbuffer[fig:einzelspalt]}

\startsectionlevel[title={Gitterkonstruktion bei
Interferenz},reference={gitterkonstruktion-bei-interferenz}]

\startitemize[packed]
\item
  S. 171 B1 und B2
\stopitemize

Drehung des Gitters: Kreise

Kreise auch mit {\bf sehr} vielen gegeneinander verdrehten Gittern

\stopsectionlevel


\startsectionlevel[title={Versuch:
Elektronenbeugungsröhre},reference={versuch-elektronenbeugungsröhre}]

\startdefinition[title={Berechnung geschwindigkeite von Elektronen in der Elektronenkanone}]

Elektronen werden im homogenen elektrischen Feld (formal Kondensator)
beschleunigt.

Im Feld besitzen sie $E_{el}=U_B \cdot e$ ($U_B$ :
Beschleunigungsspannung und $e$ : Elementarladung) an Elektrischer
Energie. Diese wird (gemäß Annahme vollständig) in {\em kinetische}
Energie umgewandelt.

\startformula 
\startalign[n=3]
\NC E_{el}                  \NC = E_{kin}                   \NC \NR
\NC e \cdot U_B             \NC = \frac{1}{2} m \cdot v^2   \NC \quad \| \cdot 2 \| \div m \NR
\NC \frac{2eU_B}{m}         \NC = v^2                       \NC \quad \| \sqrt{} \NR
\NC \sqrt{\frac{2eU_B}{m}}  \NC = v                         \NC \NR
\stopalign
 \stopformula

mit $m$ Masse Elektron und und $v$ Geschwindigkeit Elektron.

\stopdefinition

% \hrule

\startbuffer[fig:mp:beugungskanone:ringe]
\startMPpage

u := 2cm;

z0=-z3=(-2u, 0);
z1=(2u, 1u);
z2=(2u, -1u);

for a = -0.1 step 0.02 until 0.1:
	draw (x0 - a * u, 0.1 * u)--(x0 - a * u, -0.1 * u) withcolor .6white;
endfor;

path c;
c = fullcircle scaled 2.5u shifted z0;

draw c cutafter (z0 -- z1);

draw z0--z1;
draw z0--z2;

draw z0--z3;
draw z1--z3;

draw fullcircle yscaled 2u xscaled 1u shifted z3;

dotlabel("", z3);

% label.lft("R", 0.5[z1,z2]);
label.top(btex $\alpha$ etex, (-1*u, 0));
label.bot(btex $l$ etex, .5[z0,z3]);
label.rt(btex $R$ etex, .5[z1,z3]);

label.bot("Graphitfolie", (xpart z0, -0.1 * u));


\stopMPpage

\stopbuffer

\vtop{
\margintext{
	\starttabulate[|m|p|]
	\NC D \EQ Durchmesser Ring \NC\NR
	\NC R \EQ $\frac{D}{2}$ \NC\NR
	\NC l \EQ Abstand Graphit zum Leuchtschirm \NC\NR
	\stoptabulate
}
}

\placefigure[force][fig:mp:beugungskanone:ringe]{Ringe im Interferenzmuster}{
\typesetbuffer[fig:mp:beugungskanone:ringe]
}

Dazu gehört die \in{Gleichung}[formel:beugungskanone:ringe].

\startplaceformula[formel:beugungskanone:ringe]
\startformula
\tan \alpha = \frac{R}{l}
\stopformula
\stopplaceformula

\margintext{
	\starttabulate[|m|p|]
	\NC a \EQ Abstand der Spalte \NC\NR
	\NC n \EQ Ordnung \NC\NR
	\NC \lambda \EQ Wellenlänge \NC\NR
	\stoptabulate
}

Woraus sich...


\startformula
a \cdot \sin \alpha = n \cdot \lambda
\stopformula

... ableiten läst.

Eingesetzt und umgeformt sieht diese Formel wie folgt aus:

\startformula
\lambda = \frac{a \cdot \sin(\text{arctan} \frac{R}{L})}{n}
\stopformula


\startframedtask
Messbeispiele: 

\startitemize[packed]

\item Abstand des Leuchtschirems: $L = 130 mm$,
\item Durchmesser des Glaskolbens: $D = 100mm$,
\item Gitterkonstaanten: $d_1 = 123 pm; d_2 = 213 pm$

\stopitemize
\stopframedtask

\startframedtask
$U$ zu $v$
$D_1$ und $D_2$ zu $\lambda$
\stopframedtask

\startformula 
\lambda = \frac{a \cdot \sin(\text{arctan}\frac{R}{l})}{n}
 \stopformula


\stopsectionlevel


\startframedtipp
Abi aufgaben wellen: 2020 Wellen 2
\stopframedtipp
